\chapter{\IfLanguageName{dutch}{HLASM-analyse}{HLASM-analysis}}%
\label{ch:hlasmanalyse}

In dit onderzoek worden de termen "assembler" en "HLASM" gebruikt als synoniemen voor de High Level Assembler van IBM voor z/Architecture. De analyse wordt uitgevoerd op basis van de verkregen kennis vanuit \textcite{ehrman_2001} en \textcite{hlasm_language_reference}. 

\section{Inleiding}

High Level Assembler (HLASM) is een low-level programmeertaal voor IBM mainframes. Een HLASM-programma bestaat uit een reeks instructies die direct met de hardware interageren. Er zijn vier soorten instructies: commentaarinstructies, machine-instructies, assemblerinstructies en macroinstructies. 

\begin{itemize}
    \item \textbf{Commentaarinstructies}; 
        Deze instructies bieden toelichting in het programma voor de leesbaarheid en begrijpbaarheid. 
    \item \textbf{Machine-instructies}; 
        Deze instructies worden omgezet naar in machinetaalinstructies, die de cpu uitvoert wanneer het programma wordt uitgevoerd. 
    \item \textbf{Assemblerinstructies};  
        Deze instructies geven informatie aan de assembler. Ze kunnen bijvoorbeeld gebruikt worden om gegevens te genereren of een titel te geven aan de uitvoer. 
    \item \textbf{Macroinstructies}; 
        Deze instructies worden gebruikt als referentie naar een groep van instructies. De assembler vervangt die referentie automatisch naar de groep instructies. 
\end{itemize}

\newpage

De bovenvermelde instructies, met uitzondering tot commentaarinstructies, bestaan elk uit vier delen die velden/fields worden genoemd. 

\begin{itemize}
    \item The name/label field 
    \item The operartion/mnemonic field (altijd verplicht)
    \item The operand field (bij meerdere operands zijn deze gescheiden door een komma)
    \item The comment/remark field
\end{itemize}

Elk van deze velden zijn gescheiden bij één of meerdere spaties.

\begin{figure}[h!]
    \centering
    \includegraphics[width=0.75\textwidth]{asm_structuur.png}
    \caption[De opstelling van een HLASM-programma.]{De opstelling van een HLASM-programma \autocite{hlasm_language_reference}.} 
    \label{fig:instructies}
\end{figure}

\newpage

\section{De kolomgebaseerde structuur}

De Assembler verwerkt records van precies 80 bytes lang, dus dat betekent 80 tekens in één lijn code. Het aantal tekens wordt gerepresenteerd als kolommen. De instructies nemen enkel de posisties 1 tot en met 71 in beslag. Daarnaast, dient kolom 72 om mee tegeven als de volgende lijn code een voortzetting is van de huidig of een volledig nieuwe regel. Wanneer een niet-leeg teken ingevuld is, betekent dit dat het een voortzetting is. Bij een voortzetting, start de volgende regel vanaf kolom 16. De kollomen 73 tot 80 worden door de assembler genegeerd. 

De kolommen die in beslag worden genomen door de instructies (1 - 71), zijn nogmaals opgedeeld voor de verschillende velden: 

\begin{table}[h!]
    \centering
    \small
    \begin{tabular}{|p{3cm}|p{12cm}|}
        \hline
        \textbf{Kolommen} & \textbf{Inhoud} \\
        \hline
        1 - 8 & De name/label \\
        \hline
        9 & Blank \\
        \hline
        10 - 14 & De operartion/mnemonic \\
        \hline
        15 & Blank  \\
        \hline
        16 - 71 & De operands(s), gescheiden door komma's - hier mogen geen spaties voorkomen (enkel tussen quotes). Wanneer een spatie voorkomt, wordt dit gezien als het einde van de operands veld. Alles na de operands wordt dan gezien als het commentaar/remark veld. \\
        \hline
    \end{tabular}
    \caption{De kolomgebaseerde structuur van een hlasm-instructie.}
    \label{tab:kolommen}
\end{table}

\subsection{Voorbeeld}

...

\begin{lstlisting}
************************************************************************
* PROGRAMMER:  EDWARD BOSWORTH, COLUMBUS STATE UNIVERSITY              *
* SUMMARY   :  SET UP OF STANDARD LINKAGE                              *
************************************************************************
----+----1----+----2----+----3----+----4----+----5----+----6----+----7--
LAB1     CSECT                                                          
         SAVE  (14,12)                 SAVE THE CALLER'S REGISTERS     
         BALR  R12,0                   ESTABLISH                       
         USING *,R12                   ADDRESSABILITY                  
         LA    R2,SAVEAREA             ADDRESS OF MY SAVE AREA         
         ST    R2,8(,R13)              FORWARD CHAIN MINE              
         ST    R13,SAVEAREA+4          BACKWARD CHAIN CALLER'S        
         LR    R13,R2                  SET 13 FROM MY SUB CALLS        
\end{lstlisting}

\section{De control flow}

\subsection{Instructies voor Control-Flow Mapping}

